\documentclass[a4paper, 11pt]{article}

\usepackage{revy}
\usepackage[T1]{fontenc}
\usepackage{amsfonts}
\usepackage{amssymb}
\usepackage[utf8]{inputenc}
% Required for inserting images

\revyname{Matematikrevyen}
\revyyear{2025}
\version{0.1}
\eta{$3$ minutter}
\status{ Færdig}

\title{Barda}
\author{Malthe '20, Christian O '24, Louie '24, Sirius '23}

\begin{document}

\maketitle

\begin{roles}  
\role{I}[Nadja] Instruktør 
\role{S1}[Flemming] Studerende 1
\role{S2}[Alberte F] Studerende 2
\role{N}[Thais] Niels Martin
\role{In}[Ronja] Interviewer
\role{F}[Johan P] Flade
\role{N1}[Louise] Ninja 
\role{N2}[Niklas] Ninja
\end{roles}

\begin{sketch}
    





\says{In} Det er sommer på HCØ. To unge rus tager Geometri 1 for allerførste gang. De har aldrig mødt hinanden før. 

\scene En video spiller på AV: Interviews med de to deltagere i en skov, som går rundt på knæ.

\says{In} Theo er 23 år og bor på Rigshospitalets Kollegium. Han har tre frixion penne, og kan godt lide at spise koldskål. 

\says{In} Du skal snart til at have endnu et matematikkursus. Hvad ved du egentlig om geometri?

\says{S1} Der er $180^\circ$ i en trekant. 

\says{In} Okay... Hvor ved du det fra?

\says{S1} Jeg er bare rigtig god til at forestille mig ting. Og så lytter jeg meget til podcasts. 

\says{In} Emma er 19 år og bor hjemme hos forældrene i Albertslund. Hun har vundet Georg Mohr 3 gange og tager i øjeblikket også topologi. 

\says{In} Du læser matematik. Hvad kan man bruge det til?

\says{S2} Øhh det er jeg ikke helt sikker på, men jeg skal alligevel bare skrive phd og forske.

\says{In} Hvad vil du sige din stærke side er?

\says{S2} Venstre. 

\says{In} Hvordan træner man venstre?

\scene Lys på scene. Emma er allerede på scenen, Theo går akavet hen til hende. Det er meget akavet. Bandet spiller barda musik. De to studerende hilser på hinanden.

\says{S1} Hej. Jeg læser matematik. 

\says{S2} Det gør jeg også. Hvad synes du om det?

\says{S1} Det er fint nok. Har du en cykel?

\says{S2} Nej. Hvor mange point fik du til eksamen i Analyse 0?

\says{S1} Det var en mundtlig eksamen. Man får ikke point, bare en karakter. 

\says{S2} Nå. Jeg fik ellers $100$ point.

\says{N} Vær hilset! Mit navn er Nielsmester Martin. Velkommen til Geometri 1. Er I klar til et eventyr?

\says{S1} Jaaa!

\says{N} I skal nu have jeres nye navne. Theo, du er nu krigeren Theorema Egregium, en kraftfuld sætning, der binder Gauss-krumningen og lokale isometrier sammen. Emma, du er nu magikeren Lemma 2.11, et lidt ligegyldigt hjælperesultat der kan bruges til differentiere parametriserede kurver.

\says{S2} Kan vi ikke bytte? 

\says{N} Det er der slet ikke tid til – for jeres eventyr starter nu! I må skynde jer derud. Der er to veje at vælge imellem. I kan enten gå til højre eller...

\says{S1} Super, vi går til højre.

\says{S2} Nej Theorema Egregium, det er et trick. Den vej ser kortere ud lokalt, men vejen til venstre er en geodæt. Det er klart den hurtigste vej. Jeg er træt af at du trækker mig ned. JEG går til venstre!

\scene S1 og S2 splitter op. S2 går målrettet hen over scenen, mens S1 tøffer afsted og drejer rundt en masse gange. En stor flade springer ud foran S2.

\says{S2} Åh nej, det er en ikke-parametriseret overflade! 

\says{F} Waaah..

\says{S2} Hvad gør jeg?

\scene S2 vifter sin pind efter fladen.

\says{S2} Krumning! Torsion! Fundamentalformer! Et eller andet! Åh nej, der er ikke noget der virker. 

\scene Fladen ryster på hovedet, intet virker. S1 når endelig frem efter at have taget en lang omvej. Han trækker vejret forpustet og trækker sit sværd.

\says{S2} Endelig, Theorema Egregium! Der er intet det virker!

\scene S1 begynder at vifte lidt med sit sværd. S2 kigger skeptisk på.

\says{S2} Hvad laver du? Du kan ikke bare differentiere fladen!

\says{F} Waaahh... 

\scene S1 holder sit skjold oppe. Pludseligt dør fladen.

\says{F}[Døende] Waaaaaaaaaaaaaaa...... hhh

\says{S2} Egregium, det virkede! Hvad gjorde du?

\says{S1} Nå, jeg brugte bare det Mapleark jeg havde fået af en ældre studerende.
\scene Lys ned
\end{sketch}


\end{document}

